\scanpage{77}
\deleted{de la forme ${\unsure{\wp|p}}'R_2(\wp)$, $R_2$ fonction rationnelle en $\wp$, montrons que le dénominateur $Q$ de $R_2$ est une constante. Autrement, $Q(\wp)$\inserted{$=\prod(\wp\unsure{-a_i|-a})$} aurait des zéros hors de $\unsure{\mathfrak P|\mathfrak B}$, il faudrait que ces zéros fussent zéros de $\wp'$, avec une multiplicité au plus égale à celle dans $\wp'$ ; or $\wp'$ a}

\deleted{Prouvons d'abord le lemme}

\deleted{\unsure{h)|b)|la} \uline{Une fonction elliptique a au moins deux pôles non congrus mod $\mathfrak P$, ou un pôle double}}

\deleted{\deleted{\unsure{On peut le déduire du fait que} \unsure{toute}} application du tore dans la sphère a un degré topologique pair. Plus élémentairement, s'il n'y avait que deux pôles non congrus mod $\mathfrak P$, on pourrait par translation ramener \unsure{tous|donc} les pôles dans $\mathfrak P$. Si $0$ était pôle simple, \inserted{$f=\frac1z+\ldots$} $f^3+a\wp'$ serait fonction elliptique holomorphe, donc on aurait $f^3=\mathrm{c}^{\mathrm{te}}\wp'$ (b), \deleted{\illegible{1}}}

\deleted{\deleted{une \unsure{place} ; $f$ \unsure{paire} une fonction propre. $\unsure{f^2=\wp|f^2=P}$}}

% slanted
\margin{\deleted{On se sert du théorème : une fonction elliptique prend \unsure{le|un} $=$ nombre de fois sur $C/\mathfrak P$ n'importe quelle valeur complexe (finie ou infinie)}}

\deleted{$\wp'$ aura autant de zéros (dans $C/\mathfrak P$) \inserted{que de} pôles, savoir $3$, \deleted{et on vérifie aisément, en} d'autre part, comme \deleted{\unsure{d'autre}} \deleted{(b) que $\omega_1$ et $\omega_2$ sont ses pôles ; en \unsure{se donnant}} fonction elliptique \uline{impaire} n'ayant pas de pôles en $\omega_1,\omega_2,-\omega_3=\omega_1+\omega_2)$ $\wp'$ s'annule en les points ($f(\omega_i)=-f(-\omega_i)$, or $f(-\omega_i)=f(\omega_i)$ puisque $2\omega_i$ est une période). Il suit, \unsure{ainsi|ici} que \uline{$\wp'(u)$ \deleted{\unsure{en}} \inserted{a} $\omega_1,\omega_2,\omega_3$ comme pôles simples, (et pas d'autres dans $P$)}. $f'$ n'est donc pas \deleted{\unsure{un cube}} de la forme $f^3$.}

% archive sheet number: 76
